\startcomponent ctx-enlaces
\product ppchtextut


\startchapter[title={Enlaces}]

  \startsection[title=Enlaces sencillos y dobles]
    
    Empezaremos con ejemplos en los que \type{S} significa enlace sencillo y
    \type{D} enlace doble.

    En nuestro primer ejemplo representaremos cuatro plantillas en los que
    se representan enlaces sencillos y dobles, tanto en el anillo, como en
    los radicales. De paso, iremos recordando conceptos estudiados en el capítulo
    anterior:
    
    % El comando comentado numera la fórmula para poderla referenciar más tarde
    %\placeformula
    \startformula
      \startlinealignment[n=4, number=no, align={middle,middle,middle,middle}]
        % Columna 1: Estructura 1
        \NC
        \startchemical[frame=off,left=1800,width=3600,top=1900,bottom=2200]
          \chemical[SIX,B,+SR]
          \bottext{[SIX,B,+SR]}
        \stopchemical
        % Columna 2: Estructura 2
        \NC
     \startchemical[frame=off,left=1800,width=3600,top=1900,bottom=2200]          
          \chemical[SIX,SB,SR]
          \bottext{[SIX,SB,SR]}
        \stopchemical
        % Columna 5: Estructura 3
        \NC
     \startchemical[frame=off,left=2500,width=5000,top=1900,bottom=2200]          
          \chemical[SIX,SB135,DB246,SR]
          \bottext{[SIX,SB135,DB246,SR]}
        \stopchemical
        \NC
     \startchemical[frame=off,left=2500,width=5000,top=1900,bottom=2200]          
          \chemical[SIX,SB1356,DB24,DR6,SR1..5]
          \bottext{[SIX,SB1356,DB24,DR6,SR1..5]}
        \stopchemical
        \NR
      \stoplinealignment 
    \stopformula

    \startitemize[n, packed, joinedup]
    \item El ejemplo de la izquierda muestra una estructura hexagonal [\type{SIX}],
      con todos los enlaces de vértice a vértice [\type{B}]. Los radicales son
      enlaces sencillos que salen de los vértices y están acortados [\type{+SR}],
      ---recordemos que es el signo (\type{+}) el que acorta el radical al final---.
      Por último, al final de los radicales acortados situamos los \emph{átomos},
      con [\type{RZ}], en este caso números, para recordar la numeración en esta
      estructura. Todos los enlaces son sencillos.
    \item En el siguiente ejemplo seguimos con enlaces sencillos, pero ahora dejamos
      espacio, tanto para los átomos al final de los radicales, como para los átomos
      en los vértices de la estructura, con [\type{SR}], lo que acorta tanto el
      principio como el final del radical.
    \item En el tercer ejemplo nos aparecen enlaces dobles [\type{DB246}] entre los
      átomos 23, 45 y 61, debido a la numeración 246 que acompaña.
    \stopitemize

    De los tres ejemplos, infiero que [\type{SR}] dibuja radicales como enlaces
    simples entre átomos de la estructura y los que se unen a ellos, {\bf dejando
      espacio para la representación de los átomos}. De la misma forma,
    [\type{DB}] hace los mismo, {\bf pero ahora los enlaces son dobles}.
    Se ve claramente que (\type{S}) significa enlace sencillo y (\type{D}) es
    enlace doble.


    \startchemical[frame=off,left=2500,width=5000,top=1900,bottom=2500]
      \chemical[SIX,+DB246,+SB135,+DR,RZ][1,2,3,4,5,6]
      \bottext{[SIX,B,+DR,RZ] [1..6]}
    \stopchemical

    
  \stopsection

  
  \startsection[title=Enlaces sencillos orientados en el espacio]
    Empezaremos con enlaces sencillos pero que indican una orientación en el
    espacio cuando interesa esta.

    El siguiente ejemplo sirve como modelo utilizando la estructura [\type{ONE}].

    \startformula
      \startlinealignment[n=2, align={middle,middle}]
        % Columna 1: Estructura 1
        \NC
        \startchemical[frame=off, left=4000,width=8000,bottom=1600]
          \chemical[ONE,Z0,BB2,SD1,SR47,RZ1247][C,H,H,H,H]
          \bottext{[ONE,Z0,BB2,SD1,SR47,RZ1247][C,H,H,H,H]}
        \stopchemical
        % Columna 2: Estructura 2
        \NC
        \startchemical[frame=off, left=4000,width=8000,bottom=1600]
          \chemical[ONE,Z0=C,BB2,BD1,SR47,RZ1247][H,H,H,H]
          \bottext{\tfc[ONE,Z0=C,BB2,BD1,SR47,RZ1247][H,H,H,H]}
        \stopchemical
        \NR
      \stoplinealignment
    \stopformula

    \startitemize[n,packed,joinedup]
    \item En el primer ejemplo vemos una forma de representar espacialment la
      molécula de metano. Recordemos que [\type{Z0}] representa al átomo
      central, que se puede situar tanto en la lista de átomos como acompañando a
      \type{Z0}. En el código de la izquierda se sitúa en la lista de átomos, en
      cuyo caso debe estar obligatoriamente en primer lugar; en nuestro ejemplo es el
      átomo de carbono.
      Se especifican cuatro radicales como enlaces sencillos.
      Utilizaremos tres tipos diferentes para indicar la orientación espacial.
      Todos se acortan al principio y al final para dejar espacio a los átomos que
      se enlazan:
      \startitemize[packed,joinedup]
      \item Los enlaces normales que se supone que están en el plano del papel
        se reprentan como enlaces simples normales con [\type{SR}].
      \item El enlace sencillo, en negrita, que sale hacia nosotros (hacia fuera de
        la imagen), se representa mediante el código [\type{BB}].
      \item El enlace sencillo en trazos discontinuos que se esconde dentro de la
        imagen, alejándose de nosotros, se representa mediante el código [\type{SD}].
      \stopitemize      

      Por último, el código \type{ZR} sirve para indicar la orientación de los átomos
      que se unen, y sus posiciones deben relacionarse con la de los enlaces.
      La numeración debe ser 1, 2, 4 y 7, que la relacionamos con los puntos
      cardinales:
      \startitemize[packed,joinedup]
      \item {\bf 1}: Es el enlace y/o átomo que se encuentra a la derecha
        (\type{E}, \quotation{este}).
      \item {\bf 2}: El que le sigue en sentido horario
        (\type{SE}, \quotation{sureste}).
      \item {\bf 4}: El que seguiría al 3 (\type{SO}, \quotation{suroeste}).
      \item {\bf 7}: El que se encuentra arriba (\type{N}, \quotation{N}).
      \stopitemize
    \item En el segundo ejemplo encontramos dos cambios:
      \startitemize[packed,joinedup]
      \item Hemos sacado el carbono de la lista de átomos, asociándolo directamente a
        [\type{Z0}], mediante el código [\type{Z0=C}], lo que parece más intuitivo.
      \item Por otro lado, el enlace sencillo que se aleja de nosotros hacia
        dentro del papel se representa, también de manera más intuitiva como
        [\type{BD}].
      \stopitemize
    \stopitemize

%    \startchemical
%      \chemical[SIX,B,ZN] [C,C,C,C,C,C]
%    \stopchemical
%
%    \startchemical
%      \chemical[SIX,B,ZT] [1,2,3,4,5,6]
%    \stopchemical

    
  \stopsection
  
      
    \stopchapter
  \stopcomponent



  

%%% Local Variables:
%%% coding: utf-8
%%% mode: context
%%% ConTeXt-Mark-version: "IV"
%%% TeX-engine: default
%%% TeX-master: "../ctx-ppchtextut.tex"
%%% End:
